%; whizzy chapter -dvi
% -initex iniptex -latex platex -format platex -bibtex jbibtex -fmt fmt
% 以上 whizzytex を使用する場合の設定。
 
%     Tokyo Debian Meeting resources
%     Copyright (C) 2012 Junichi Uekawa
%     Copyright (C) 2011, 2015 Nobuhiro Iwamatsu

%     This program is free software; you can redistribute it and/or modify
%     it under the terms of the GNU General Public License as published by
%     the Free Software Foundation; either version 2 of the License, or
%     (at your option) any later version.

%     This program is distributed in the hope that it will be useful,
%     but WITHOUT ANY WARRANTY; without even the implied warranty of
%     MERCHANTABILITY or FITNESS FOR A PARTICULAR PURPOSE.  See the
%     GNU General Public License for more details.

%     You should have received a copy of the GNU General Public License
%     along with this program; if not, write to the Free Software
%     Foundation, Inc., 51 Franklin St, Fifth Floor, Boston, MA  02110-1301 USA

%  preview (shell-command (concat "evince " (replace-regexp-in-string "tex$" "pdf"(buffer-file-name)) "&"))

%%ここからヘッダ開始。

\documentclass[mingoth,a4paper]{jsarticle}
\usepackage{monthlyreport}
% 日付を定義する、毎月変わります。
\newcommand{\debmtgyear}{2016}
\newcommand{\debmtgmonth}{11}
\newcommand{\debmtgdate}{19}
% started from zero:
% (let ((year 2013) (month 7)) (+ (* (- year 2005) 12) month -1))
\newcommand{\debmtgnumber}{146}

\begin{document}

\begin{titlepage}
\thispagestyle{empty}
% タイトルページ:編集必要な部分は最初のマクロに飛ばすこと

\vspace*{-2cm}
第\debmtgnumber{}回 東京エリア Debian 勉強会資料\\
\hspace*{-2cm}
\includegraphics{image2012-natsu/dotdeb.pdf}\\
\hfill{}\debmtgyear{}年\debmtgmonth{}月\debmtgdate{}日

% ここはアップデートすること
% 全角文字にしないとフォントのサイズが合わないので注意
\rotatebox{10}{\fontsize{30}{30} {\gt 特集 ：dh\_strip\_nondeterminism}}\\

\vspace*{-2cm}
\hfill{}\includegraphics[height=6cm]{image200502/openlogo-nd.eps}
\end{titlepage}

\newpage

\begin{minipage}[b]{0.2\hsize}
 \definecolor{titleback}{gray}{0.9}
 \colorbox{titleback}{\rotatebox{90}{\fontsize{80}{80} {\gt デビアン勉強会} }}
\end{minipage}
\begin{minipage}[b]{0.8\hsize}
\hrule
\vspace{2mm}
\hrule
\begin{multicols}{2}
\tableofcontents
\end{multicols}
\vspace{2mm}
\hrule
\end{minipage}

\dancersection{最近のDebian関連のミーティング報告}{杉本 典充}

\subsection{第144回東京エリアDebian勉強会}

2016年10月15日(土)に第144回東京エリアDebian勉強会を開催しました。会場は銀座にある朝日ネットさんをお借りして行いました。参加者は6名でした。発表は、dictossさんによる「preseedでDebianを自動インストールをしてみよう」でした。

昨今VPSサービスやPCで仮想マシンを扱うことが増えてきたことを受け、Debianを自動インストールするpreseed機能の解説と自動インストールの実演を行いました。

勉強会終了後は参加者で懇親を深めました。

\subsection{OSC 2016 Tokyo/Fall(第145回東京エリアDebian勉強会)}

2016年11月5日(土)、11月6日(日)に開催されたOSC 2016 Tokyo/Fallへ、東京エリアDebian勉強会/DebianJPとして出展しました。

展示ブースではDebian GNU/Linux及びDebian GNU/kFreeBSDをインストールしたマシンの展示、Debian Projectと本勉強会の広報活動を行いました。また、OSC参加者と交流を深めました。

セミナーでは第145回東京エリアDebian勉強会として岩松さんによる「Debian updates」を発表し、参加者からの疑問や質問に回答しました。

\dancersection{事前課題}{杉本 典充}

今回の事前課題は以下です:
\begin{enumerate}
  \item Hack Timeは何をしますか。
\end{enumerate}
この課題に対して提出いただいた内容は以下です。
\begin{multicols}{2}
{\small
\begin{prework}{ user }
  \begin{enumerate}
  \item xx
  \end{enumerate}
\end{prework}

}
\end{multicols}

%\dancersection{Debian Trivia Quiz}{杉本 典充}
%
%Debianの昨今の話題についてのQuizです。
%
%今回の出題範囲は\url{debian-devel-announce@lists.debian.org} や \url{debian-news@lists.debian.org}に投稿された内容などからです。
%
%\begin{multicols}{2}
%%; whizzy-master ../debianmeetingresume201311.tex
% 以上の設定をしているため、このファイルで M-x whizzytex すると、whizzytexが利用できます。
%

\santaku
{q}
{A}
{B}
{C}
{B}
{desc}

%\end{multicols}


% % (query-replace-regexp "<.*?>" "")
% % (query-replace-regexp "^[	 ]\+" "")

%-------------------------------------------------------------------------------
\dancersection{dh\_strip\_nondeterminism について}{吉野 与志仁}
%-------------------------------------------------------------------------------
\subsection{はじめに}

2015 年に debhelper の \verb|dh(1)| シーケンスに組み込まれた
\verb|dh_strip_nondeterminism(1)| ですが、これが
具体的に何をやっているのか調べてみました。調べたバージョンは
dh-strip-nondeterminism パッケージの 0.028-1 です。

なお、このコマンドの目的は reproducible builds（再現可能なビルド）ですが、詳細は
\url{https://wiki.debian.org/ReproducibleBuilds} や
\url{https://reproducible-builds.org/} などを見ると良いでしょう。

再現可能なビルドを実現するため、このコマンドはビルドツールが生成した
データに埋め込まれた時刻などを指定した値に書き換えます。実際に見ていきま
しょう。

\subsection{具体例}
\subsubsection{ar}
*.a のうち \verb|ar archive| のもの

\begin{commandline}
$ file /usr/lib/x86_64-linux-gnu/libglib-2.0.a
/usr/lib/x86_64-linux-gnu/libglib-2.0.a: current ar archive
\end{commandline}

\verb|dh_strip_nondeterminism(1)| は、パッケージの debian/changelog
ファイルの中に書いてある時刻を使うようになっています。
この .a ファイルを含むパッケージの debian/changelog ファイルの時刻を見てみます。
\begin{commandline}
$ dpkg -S /usr/lib/x86_64-linux-gnu/libglib-2.0.a
libglib2.0-dev: /usr/lib/x86_64-linux-gnu/libglib-2.0.a
$ dpkg-parsechangelog -l /usr/share/doc/libglib2.0-dev/changelog.Debian.gz
Source: glib2.0
Version: 2.50.2-1
Distribution: unstable
Urgency: medium
Maintainer: Michael Biebl <biebl@debian.org>
Timestamp: 1478561825
Date: Tue, 08 Nov 2016 00:37:05 +0100
Changes:
 glib2.0 (2.50.2-1) unstable; urgency=medium
 .
   * New upstream release.
   * Track stable releases in debian/watch.
\end{commandline}
この \verb|Date:|で始まる行（もしくは\verb|Timestamp:|で始まる行）の値が使われます。

ファイルの中を見てみます。
\begin{commandline}
$ env TZ=UTC 7z l /usr/lib/x86_64-linux-gnu/libglib-2.0.a | head -n 20

7-Zip [64] 16.02 : Copyright (c) 1999-2016 Igor Pavlov : 2016-05-21
p7zip Version 16.02 (locale=ja_JP.UTF-8,Utf16=on,HugeFiles=on,64 bits,4 CPUs Intel(R) Core(TM) i5-4250U CPU @ 1.30GHz (40651),ASM,AES-NI)

Scanning the drive for archives:
1 file, 1973536 bytes (1928 KiB)

Listing archive: /usr/lib/x86_64-linux-gnu/libglib-2.0.a

--
Path = /usr/lib/x86_64-linux-gnu/libglib-2.0.a
Type = Ar
Physical Size = 1973536
SubType = a

   Date      Time    Attr         Size   Compressed  Name
------------------- ----- ------------ ------------  ------------------------
2016-11-07 23:37:05 .....        84550        84550  1.txt
2016-11-07 23:37:05 .....         4584         4584  libglib_2_0_la-gallocator.o
2016-11-07 23:37:05 .....         6352         6352  libglib_2_0_la-gcache.o
$ env TZ=UTC ar tv /usr/lib/x86_64-linux-gnu/libglib-2.0.a | head -n 4
rw-r--r-- 0/0   4584 Nov  7 23:37 2016 libglib_2_0_la-gallocator.o
rw-r--r-- 0/0   6352 Nov  7 23:37 2016 libglib_2_0_la-gcache.o
rw-r--r-- 0/0   5960 Nov  7 23:37 2016 libglib_2_0_la-gcompletion.o
rw-r--r-- 0/0  11464 Nov  7 23:37 2016 libglib_2_0_la-grel.o
\end{commandline}

固められた各ファイルの
\begin{itemize}
 \item タイムスタンプ (mtime) を指定した時刻に書き換えています。
 \item 所有者 (owner) を 0 に上書きしています。
 \item グループ (group) を 0 に上書きしています。
 \item パーミッション (mode) を 755 か 644 のいずれかに揃えています。
\end{itemize}

\subsubsection{zip}

*.zip, *.pk3, *.epub, *.whl, *.xpi, *.htb, *.zhfst のうち \verb|Zip archive data|
または \verb|EPUB document| のもの

\begin{commandline}
$ file /usr/share/go-1.7/src/archive/zip/testdata/symlink.zip
/usr/share/go-1.7/src/archive/zip/testdata/symlink.zip: Zip archive data, at least v1.0 to extract
$ file /usr/share/debian-reference/debian-reference.ja.epub
/usr/share/debian-reference/debian-reference.ja.epub: Zip archive data, at least v1.0 to extract
\end{commandline}

調べたパッケージの時刻はそれぞれ
\begin{commandline}
$ dpkg -S /usr/share/go-1.7/src/archive/zip/testdata/symlink.zip
golang-1.7-src: /usr/share/go-1.7/src/archive/zip/testdata/symlink.zip
$ dpkg-parsechangelog -l /usr/share/doc/golang-1.7-src/changelog.Debian.gz | grep ^Date:
Date: Thu, 20 Oct 2016 09:10:47 +1300
\end{commandline}

\begin{commandline}
$ dpkg -S /usr/share/debian-reference/debian-reference.ja.epub
debian-reference-ja: /usr/share/debian-reference/debian-reference.ja.epub
$ dpkg-parsechangelog -l /usr/share/doc/debian-reference-ja/changelog.gz | grep ^Date:
Date: Mon, 17 Oct 2016 22:28:00 +0900
\end{commandline}
です。

中身を見てみます。

\begin{commandline}
$ env TZ=UTC 7z l /usr/share/go-1.7/src/archive/zip/testdata/symlink.zip | tail -n 5
   Date      Time    Attr         Size   Compressed  Name
------------------- ----- ------------ ------------  ------------------------
2016-10-19 20:10:47 .....            9            9  symlink
------------------- ----- ------------ ------------  ------------------------
2016-10-19 20:10:47                  9            9  1 files
\end{commandline}

\begin{commandline}
$ zipinfo -v /usr/share/go-1.7/src/archive/zip/testdata/symlink.zip | tail -n 36
Central directory entry #1:
---------------------------

  symlink

  offset of local header from start of archive:   0
                                                  (0000000000000000h) bytes
  file system or operating system of origin:      Unix
  version of encoding software:                   3.0
  minimum file system compatibility required:     MS-DOS, OS/2 or NT FAT
  minimum software version required to extract:   1.0
  compression method:                             none (stored)
  file security status:                           not encrypted
  extended local header:                          no
  file last modified on (DOS date/time):          2016 Oct 19 20:10:46
  file last modified on (UT extra field modtime): 2016 Oct 20 05:10:47 local
  file last modified on (UT extra field modtime): 2016 Oct 19 20:10:47 UTC
  32-bit CRC value (hex):                         8e9efad1
  compressed size:                                9 bytes
  uncompressed size:                              9 bytes
  length of filename:                             7 characters
  length of extra field:                          24 bytes
  length of file comment:                         0 characters
  disk number on which file begins:               disk 1
  apparent file type:                             binary
  Unix file attributes (100755 octal):            -rwxr-xr-x
  MS-DOS file attributes (00 hex):                none

  The central-directory extra field contains:
  - A subfield with ID 0x5455 (universal time) and 5 data bytes.
    The local extra field has UTC/GMT modification/access times.
  - A subfield with ID 0x7875 (Unix UID/GID (any size)) and 11 data bytes:
    01 04 00 00 00 00 04 00 00 00 00.

  There is no file comment.

\end{commandline}

\begin{commandline}
$ env TZ=JST-9 7z l /usr/share/debian-reference/debian-reference.ja.epub | head -n 24 | tail -n 10
   Date      Time    Attr         Size   Compressed  Name
------------------- ----- ------------ ------------  ------------------------
2016-10-17 22:28:00 D....            0            0  META-INF
2016-10-17 22:28:00 .....          255          175  META-INF/container.xml
2016-10-17 22:28:00 D....            0            0  OEBPS
2016-10-17 22:28:00 .....        10451         3596  OEBPS/apa.html
2016-10-17 22:28:00 .....       123568        16458  OEBPS/bk01-toc.html
2016-10-17 22:28:00 .....       279742        45365  OEBPS/ch01.html
2016-10-17 22:28:00 .....       306201        49878  OEBPS/ch02.html
2016-10-17 22:28:00 .....        86922        15031  OEBPS/ch03.html
$ env TZ=JST-9 7z l /usr/share/debian-reference/debian-reference.ja.epub | tail -n 4
2016-10-17 22:28:00 .....        96004        19619  OEBPS/toc.ncx
2016-10-17 22:28:00 .....           20           20  mimetype
------------------- ----- ------------ ------------  ------------------------
2016-10-17 22:28:00            2593315       425235  23 files, 2 folders
\end{commandline}
\begin{commandline}
$ env TZ=UTC zipinfo -v /usr/share/debian-reference/debian-reference.ja.epub OEBPS/ch01.html | tail -n 12
  apparent file type:                             text
  Unix file attributes (100644 octal):            -rw-r--r--
  MS-DOS file attributes (00 hex):                none

  The central-directory extra field contains:
  - A subfield with ID 0x5455 (universal time) and 5 data bytes.
    The local extra field has UTC/GMT modification/access times.
  - A subfield with ID 0x7875 (Unix UID/GID (any size)) and 11 data bytes:
    01 04 00 00 00 00 04 00 00 00 00.

  There is no file comment.

\end{commandline}

固められた各ファイルの
\begin{itemize}
 \item 並びを名前順に直しています。
 \item DOS 時刻のフィールドを指定した時刻に書き換えています。しかし、精
       度が2秒しかないようです。
 \item 属性を 755 か 644 のいずれかに揃えています。
 \item central directory header（と local header）の
       \begin{itemize}
	\item ID 0x5455 (\verb|UT| universal time)のフィールドに時刻が入っているの
	      で、指定した時刻に書き換えています。
	\item ID 0x7875 (\verb|ux| Unix UID/GID)のフィールドにあるUIDとGIDを0
	      に上書きしています。
       \end{itemize}
\end{itemize}

\subsubsection{jar}
*.jar, *.war, *.hpi, *.apk のうち \verb|Java archive| または \verb|Zip archive| のもの

\begin{commandline}
$ file ./usr/share/java/commons-lang3.jar
./usr/share/java/commons-lang3.jar: Zip archive data, at least v1.0 to extract
\end{commandline}

\begin{commandline}
$ dpkg-parsechangelog -l ./usr/share/doc/libcommons-lang3-java/changelog.Debian.gz | grep ^Date:
Date: Thu, 20 Oct 2016 19:08:15 +0200
\end{commandline}

中を見てみます。
\begin{commandline}
$ env TZ=UTC 7z l ./usr/share/java/commons-lang3.jar | head -n 24 | tail -n 10
   Date      Time    Attr         Size   Compressed  Name
------------------- ----- ------------ ------------  ------------------------
2016-10-20 17:08:14 D....            0            0  META-INF
2016-10-20 17:08:14 .....         1844          732  META-INF/MANIFEST.MF
2016-10-20 17:08:14 .....        11358         3949  META-INF/LICENSE.txt
2016-10-20 17:08:14 .....          301          187  META-INF/NOTICE.txt
2016-10-20 17:08:14 D....            0            0  META-INF/maven
2016-10-20 17:08:14 D....            0            0  META-INF/maven/org.apache.commons
2016-10-20 17:08:14 D....            0            0  META-INF/maven/org.apache.commons/commons-lang3
2016-10-20 17:08:14 .....           91           83  META-INF/maven/org.apache.commons/commons-lang3/pom.properties
$ env TZ=UTC 7z l ./usr/share/java/commons-lang3.jar | tail -n 4
2016-10-20 17:08:14 .....         3231         1289  org/apache/commons/lang3/tuple/Pair.class
2016-10-20 17:08:14 .....         3190         1250  org/apache/commons/lang3/tuple/Triple.class
------------------- ----- ------------ ------------  ------------------------
2016-10-20 17:08:14             986869       404745  265 files, 19 folders
\end{commandline}

\begin{commandline}
$ unzip -p ./usr/share/java/commons-lang3-3.5.jar META-INF/MANIFEST.MF
Manifest-Version: 1.0
Bundle-License: http://www.apache.org/licenses/LICENSE-2.0.txt
Bundle-SymbolicName: org.apache.commons.lang3
Archiver-Version: Plexus Archiver
Implementation-Vendor-Id: org.apache
Specification-Title: Apache Commons Lang
Bundle-DocURL: http://commons.apache.org/proper/commons-lang/
Include-Resource: META-INF/LICENSE.txt=LICENSE.txt,META-INF/NOTICE.txt
 =NOTICE.txt
Require-Capability: osgi.ee;filter:="(&(osgi.ee=JavaSE)(version=1.5))"
Export-Package: org.apache.commons.lang3;version="3.5",org.apache.comm
 ons.lang3.builder;version="3.5",org.apache.commons.lang3.concurrent;v
 ersion="3.5",org.apache.commons.lang3.event;version="3.5",org.apache.
 commons.lang3.exception;version="3.5",org.apache.commons.lang3.math;v
 ersion="3.5",org.apache.commons.lang3.mutable;version="3.5",org.apach
 e.commons.lang3.reflect;version="3.5",org.apache.commons.lang3.text;v
 ersion="3.5",org.apache.commons.lang3.text.translate;version="3.5",or
 g.apache.commons.lang3.time;version="3.5",org.apache.commons.lang3.tu
 ple;version="3.5"
Bundle-Name: Apache Commons Lang
Implementation-Title: Apache Commons Lang
Bundle-Description: Apache Commons Lang, a package of Java utility cla
 sses for the  classes that are in java.lang's hierarchy, or are consi
 dered to be so  standard as to justify existence in java.lang.
Implementation-Version: 3.5
Specification-Vendor: The Apache Software Foundation
Bundle-ManifestVersion: 2
Bundle-Vendor: The Apache Software Foundation
Tool: Bnd-2.4.1.201608301338
Implementation-Vendor: The Apache Software Foundation
Bundle-Version: 3.5.0
X-Compile-Target-JDK: 1.5
Implementation-Build: ${scmBranch}@r${buildNumber}; 2016-10-20 17:08:1
 5+0000
X-Compile-Source-JDK: 1.5
Created-By: Apache Maven Bundle Plugin
Build-Jdk: 1.8.0_102
Specification-Version: 3.5

\end{commandline}

jar は zip ファイルなので、*.zip と同じ変更も加えています。

さらに、
\begin{itemize}
 \item 各ファイルの並びを名前順に直しています。ただ \verb|META-INF/| と
       \verb|META-INF/MANIFEST.MF| は先頭にしているようです。
 \item META-INF/MANIFEST.MF から
       \begin{itemize}
	\item \verb|Bnd-LastModified:| で始まる生成時刻が書かれた行を削除しています。
	\item \verb|Built-By:| で始まるビルドしたユーザ名が書かれた行を
	      削除しています。
	\item コンパイラやビルドツールのバージョン番号は残しているようで
	      す。コンパイラのバグ修正等を想定しているのかもしれません。
       \end{itemize}
 \item *.properties も変更しています（後述）。
 \item javadoc な *.html があったら変更するらしいです（後述）。
 \item *.jar があったら再帰的に変更するらしいです。
 \item META-INF/*.SF を含んだ jar ファイルは処理の対象外にしています
       (Bug\#807669)。署名付き jar ファイルに含まれるようですが、普通は
       ビルド時に自動署名しないだろうという想定なのでしょう。
\end{itemize}

\subsubsection{javadoc}
*.html のうち \verb|<!-- Generated by javadoc| があるもの

\begin{commandline}
$ grep \<html ./usr/share/doc/libcommons-lang3-java/api/org/apache/commons/lang3/StringUtils.html
<html>
$ grep '<!-- Generated by javadoc' ./usr/share/doc/libcommons-lang3-java/api/org/apache/commons/lang3/StringUtils.html
<!-- Generated by javadoc -->
\end{commandline}

ファイル内の
\begin{itemize}
 \item 生成時の環境の言語に基づいて書かれたhtml要素のlang属性
       (\verb|<html lang=|) を削除しています。
 \item \verb|<!-- Generated by javadoc| 行から他の文字列（生成時刻や生成ツールのバー
       ジョン番号）をすべて削除しています。ドキュメントの生成ツールは多
       少バージョンが変わっても生成結果は変わらない想定なのでしょう。
\end{itemize}

\subsubsection{javaproperties}
*.properties のうち Java 系のビルドツールで自動生成されたように見える、
\verb|#Generated by Apache Maven| などを含むもの

\begin{commandline}
$ unzip -p ./usr/share/java/commons-lang3-3.5.jar META-INF/maven/org.apache.commons/commons-lang3/pom.properties
#Generated by Apache Maven
version=3.5
groupId=org.apache.commons
artifactId=commons-lang3
\end{commandline}

\verb|#| で始まる自動生成時刻が書かれた行を削除しています。

\subsubsection{png}
*.png のうち \verb|PNG image data| のもの

\begin{commandline}
$ file /usr/share/doc/debian-handbook/html/ja-JP/images/developers-map.png
/usr/share/doc/debian-handbook/html/ja-JP/images/developers-map.png: PNG image data, 750 x 450, 8-bit/color RGB, non-interlaced
$ file /usr/share/emacs/24.5/etc/images/splash.png
/usr/share/emacs/24.5/etc/images/splash.png: PNG image data, 275 x 188, 8-bit/color RGBA, non-interlaced
\end{commandline}

\begin{commandline}
$ dpkg -S /usr/share/doc/debian-handbook/html/ja-JP/images/developers-map.png
debian-handbook: /usr/share/doc/debian-handbook/html/ja-JP/images/developers-map.png
$ dpkg-parsechangelog -l /usr/share/doc/debian-handbook/changelog.gz | grep ^Date:
Date: Thu, 22 Sep 2016 16:09:44 +0200
\end{commandline}

\begin{commandline}
$ dpkg -S /usr/share/emacs/24.5/etc/images/splash.png
emacs24-common: /usr/share/emacs/24.5/etc/images/splash.png
$ dpkg-parsechangelog -l /usr/share/doc/emacs24-common/changelog.Debian.gz | grep ^Date:
Date: Mon, 05 Sep 2016 15:05:00 -0500
\end{commandline}

中を見てみます。
\begin{commandline}
$ hd /usr/share/doc/debian-handbook/html/ja-JP/images/developers-map.png | grep tIME
00000070  00 00 07 74 49 4d 45 07  e0 09 16 0e 09 2c 4b 76  |...tIME......,Kv|
#               ^^
#              長さ
#                              ^^^^^^ ^^ ^^ ^^ ^^ ^^
#                               2016  09 22 14 09 44
$ strings -a /usr/share/doc/debian-handbook/html/ja-JP/images/developers-map.png | grep -A1 '[tiz]EXtdate:'
%tEXtdate:create
2016-09-22T14:09:44-00:00
%tEXtdate:modify
2016-09-22T14:09:44-00:00
$ exiftool /usr/share/doc/debian-handbook/html/ja-JP/images/developers-map.png | tail -n 5
Modify Date                     : 2016:09:22 14:09:44
Datecreate                      : 2016-09-22T14:09:44-00:00
Datemodify                      : 2016-09-22T14:09:44-00:00
Image Size                      : 750x450
Megapixels                      : 0.338
\end{commandline}

変更時刻 (\verb|tIME|)、\verb|date:| の値を指定した時刻に書き換えています。

\begin{commandline}
$ strings -a /usr/share/emacs/24.5/etc/images/splash.png | grep -A1 '[tiz]EXtCreation Time'
'tEXtCreation Time
2016-09-05T20:05:00-00:00
$ exiftool /usr/share/emacs/24.5/etc/images/splash.png | tail -n 4
Description                     : GNU Emacs splash image
Creation Time                   : 2016-09-05T20:05:00-00:00
Image Size                      : 275x188
Megapixels                      : 0.052
\end{commandline}

作成時刻 (\verb|Creation Time|) を指定した時刻に書き換えています。

自動生成しないアートワーク等もあると思いますが、どちらか区別が付かないから全部書き
換えているのでしょうか。

\subsubsection{gettext}
*.mo, *.gmo のうち \verb|GNU message catalog| のもの

\begin{commandline}
$ file /usr/share/locale/ja/LC_MESSAGES/grub.mo
/usr/share/locale/ja/LC_MESSAGES/grub.mo: GNU message catalog (little endian), revision 0.0, 233 messages
$ file /usr/share/locale/ja/LC_MESSAGES/apt.mo
/usr/share/locale/ja/LC_MESSAGES/apt.mo: GNU message catalog (little endian), revision 0.0, 367 messages
\end{commandline}

各パッケージの時刻は
\begin{commandline}
$ dpkg -S /usr/share/locale/ja/LC_MESSAGES/grub.mo
grub-common: /usr/share/locale/ja/LC_MESSAGES/grub.mo
$ dpkg-parsechangelog -l /usr/share/doc/grub-common/changelog.Debian.gz | grep ^Date:
Date: Tue, 01 Nov 2016 11:10:52 +0000
\end{commandline}

\begin{commandline}
$ dpkg -S /usr/share/locale/ja/LC_MESSAGES/apt.mo
apt: /usr/share/locale/ja/LC_MESSAGES/apt.mo
$ dpkg-parsechangelog -l /usr/share/doc/apt/changelog.gz | grep ^Date:
Date: Tue, 04 Oct 2016 19:43:35 +0200
\end{commandline}
です。

中を見てみると
\begin{commandline}
$ grep -a POT-Creation-Date: /usr/share/locale/ja/LC_MESSAGES/grub.mo
POT-Creation-Date: 2016-11-01 11:10+0000
\end{commandline}

\verb|POT-Creation-Date:| で始まる行に生成時刻が入っているので、指定した
時刻に書き換えています。

ただ、指定した時刻より新しい場合のみ上書きしています。

\begin{commandline}
$ grep -a POT-Creation-Date: /usr/share/locale/ja/LC_MESSAGES/apt.mo
POT-Creation-Date: 2016-08-30 22:20+0200
\end{commandline}

指定した時刻より古いので書き換えていません。

\subsubsection{pearregistry}

*.reg のうち先頭が \verb|a:| のもの

\begin{commandline}
$ file ./usr/share/php/.registry/services_weather.reg
./usr/share/php/.registry/services_weather.reg: ASCII text, with very long lines
$ hd ./usr/share/php/.registry/services_weather.reg | head -n 4
00000000  61 3a 32 32 3a 7b 73 3a  37 3a 22 61 74 74 72 69  |a:22:{s:7:"attri|
00000010  62 73 22 3b 61 3a 36 3a  7b 73 3a 31 35 3a 22 70  |bs";a:6:{s:15:"p|
00000020  61 63 6b 61 67 65 72 76  65 72 73 69 6f 6e 22 3b  |ackagerversion";|
00000030  73 3a 35 3a 22 31 2e 39  2e 34 22 3b 73 3a 37 3a  |s:5:"1.9.4";s:7:|
\end{commandline}

\begin{commandline}
$ dpkg-parsechangelog -l ./usr/share/doc/php-services-weather/changelog.Debian.gz | grep -e ^Date: -e ^Timestamp:
Timestamp: 1460011442
Date: Thu, 07 Apr 2016 08:44:02 +0200
\end{commandline}

\begin{commandline}
$ hd ./usr/share/php/.registry/services_weather.reg | tail -n 8
00003580  65 22 3b 73 3a 33 3a 22  65 72 75 22 3b 73 3a 34  |e";s:3:"eru";s:4|
00003590  3a 22 72 6f 6c 65 22 3b  73 3a 34 3a 22 6c 65 61  |:"role";s:4:"lea|
000035a0  64 22 3b 7d 7d 7d 73 3a  31 30 3a 22 78 73 64 76  |d";}}}s:10:"xsdv|
000035b0  65 72 73 69 6f 6e 22 3b  73 3a 33 3a 22 32 2e 30  |ersion";s:3:"2.0|
000035c0  22 3b 73 3a 31 33 3a 22  5f 6c 61 73 74 6d 6f 64  |";s:13:"_lastmod|
000035d0  69 66 69 65 64 22 3b 69  3a 31 34 36 30 30 31 31  |ified";i:1460011|
000035e0  34 34 32 3b 7d                                    |442;}|
000035e5
\end{commandline}

\verb|_lastmodified| に生成時刻が書かれているので、指定した時刻に書き換えています。

\subsubsection{gzip}

*.gz, *.dz のうち \verb|gzip compressed data| のもの


\begin{itemize}
 \item ファイル名 (FNAME) の削除
 \item ヘッダCRC (FHCRC) の削除
 \item ヘッダ内の時刻 (mtime) が空でなく新しい場合は書き換え
\end{itemize}
をするらしいんですが、現物を探せませんでした。というのも、Debianパッケージ
内によくある*.gzファイル（今までたくさん見てきたchangelogファイルやmanペー
ジなど）
は\verb|dh_compress(1)|がファイル名や時刻を保存しないオプションで圧縮し
ています。しかも\verb|dh(1)|シーケンスでは
\verb|dh_strip_nondeterminism(1)| の後が \verb|dh_compress(1)|なので書き
換え対象になっていません。

\begin{commandline}
$ grep -C2 dh_strip_nondeterminism /usr/bin/dh
	dh_installwm
	dh_installxfonts
	dh_strip_nondeterminism
	dh_compress
	dh_fixperms
\end{commandline}

\subsection{まとめ}

\verb|dh_strip_nondeterminism(1)| は、再現可能なビルドにしていくに当たり、
ビルド結果（バイナリなど）を使う際には通常不要な生成時刻などの情報を決まった
値に書き換えています。

ビルドツール等を開発している方は、これに頼らず、そもそもビルド時にこういっ
た情報を埋め込まないようにしたほうがよいでしょう。

普通のパッケージメンテナの方は、このコマンドはこのような書き換えをしているので
（manページには\verb|good guesses|していると書いてありますが）バグっている
ものがあったら報告しましょう。

reproducible builds に興味がある方は、もっと書き換えたほうがいいかもしれ
ないものがあったらバグ報告したほうがよいかもしれません。

\subsection{参考文献}

dh-strip-nondeterminism 0.028-1, p7zip-full, unzip, zip, libarchive-zip-perl,
libimage-exiftool-perl 各パッケージ

%
% 冊子にするために、4の倍数にする必要がある。
% そのための調整
%\dancersection{メモ}{}
\mbox{}\newpage


\vspace*{15cm}
\hrule
\vspace{2mm}
\includegraphics[width=2cm]{image200502/openlogo-nd.eps}
\noindent \Large \bf Debian 勉強会資料\\
\noindent \normalfont \debmtgyear{}年\debmtgmonth{}月\debmtgdate{}日 \hspace{5mm}  初版第1刷発行\\
\noindent \normalfont 東京エリア Debian 勉強会 （編集・印刷・発行）\\
\hrule

\end{document}
